%*******************************************************************************
%****************************** Chapter 8 *************************************
%*******************************************************************************

\chapter{Conclusions}
\label{ch:conclusions}

\ifpdf
    \graphicspath{{Chapter8/Figs/Raster/}{Chapter8/Figs/PDF/}{Chapter8/Figs/}}
\else
    \graphicspath{{Chapter8/Figs/Vector/}{Chapter8/Figs/}}
\fi


%********************************** 8.1 **************************************
\section{Summary}
\label{sec:concl-summary}

% TODO: summarize the work in roughly one page: the problem, what was built, and what was
% measured. Written last, from Chapters 5 to 7, and must not introduce any result not already
% reported there.


%********************************** 8.2 **************************************
\section{Conclusions}
\label{sec:concl-conclusions}

% TODO: state one conclusion per objective from Section~\ref{sec:objectives}, each supported by a
% specific measured result. Avoid claims the validation does not reach: internal testing
% establishes that the prototype works, not that it improves learning outcomes, which was never
% measured.


%********************************** 8.3 **************************************
\section{Recommendations}
\label{sec:recommendations}

% TODO: recommendations for future work. Those already declared as limitations in
% Section~\ref{subsec:limitations} belong here:
%   - user-feedback studies and classroom beta testing with CHE 3214L and CHE 4117L cohorts
%   - long-term reliability and multi-semester deployment
%   - formal adoption as a scheduled laboratory experiment
%   - additional control schemes beyond PI speed and on-off level, now that the platform exists
%   - replication: building further units so that groups can work in parallel